\documentclass[a4paper]{article}

\usepackage{graphicx}
\usepackage{amsmath,amssymb}
\usepackage{minted}
\usepackage{multirow}
\usepackage{float}
\usepackage{todonotes}
\usepackage{forest}
\usepackage{xstring}
\usepackage{xcolor}
\usepackage[most]{tcolorbox}
\usepackage{xparse}
\usepackage{natbib}

\usetikzlibrary{graphs,graphdrawing}
\usegdlibrary{layered}

% for external graphviz
\input{/home/donkarlo/Dropbox/repo/nd_ai_project/src/nd_ai/knoweldge_representation/language/formal/computer/programming/tex/latex/code_clips/external_graphviz.tex}

% for tags
\input{/home/donkarlo/Dropbox/repo/nd_ai_project/src/nd_ai/knoweldge_representation/language/formal/computer/programming/tex/latex/parts/control_sequence/kind/semtag/semtag.tex}

% for links
\input{/home/donkarlo/Dropbox/repo/nd_ai_project/src/nd_ai/knoweldge_representation/language/formal/computer/programming/tex/latex/code_clips/seehere.tex}

\title{Glossematics}
\date{}
\author{Mohammad Rahmani}

\begin{document}
\maketitle
    Glossematics defines the glosseme as the smallest irreducible unit of both the content and expression planes of language; in the expression plane, the glosseme is nearly identical to the phoneme. In the content plane, it is the smallest unit of meaning which underlies a concept. A ewe, for example, consists of the taxemes sheep and female which may eventually be divided into even smaller units – glossemes – of meaning. The analysis is gradually expanded to the study of functions, more commonly known as dependencies, between elements on the level of discourse (which is called process), and between meaning and form in the linguistic system  \seeherefooter{https://en.wikipedia.org/wiki/Glossematics}.
    \bibliographystyle{plainnat}
    \bibliography{/home/donkarlo/Dropbox/repo/nd_shared_research/bibliography/bibliography.bib}
\end{document}